\scanpage{101}

\margin{\tikz[baseline]{\draw[line width=.7pt] (0,0)--(0,1.6);}}
\uline{Corollaire} 1 $u\in L^N(E,F)$ si et seulement si $|u|\in L^N(E,F)$, et

\[
N(u)=N(|u|)
\]

$u\in L_N(E,F)$ si et seulement si $|u|\in L_N(E,F)$.

\deleted{\uline{Corollaire 2} $N(u)$ est fonction croissante sur l'ensemble des op. hermitiens positifs dans $E$. En effet, $\unsure{(\rho_i(u))|(\rho_i u)}$ est fonction croissante de $u$.}

\uline{Corollaire 2} Si $u=h+ik$, $h$ et $k$ hermitiens, alors

\[
u\in L_N(E)\Longleftrightarrow h,k\in L_N(E),\qquad
u\in L^N(E)\Longleftrightarrow h,k\in L^N(E).
\]

\uline{Remarque} Soit $p(u)$ une norme sur $E'\otimes F$ satisfaisant aux conditions\,: a)

\[
\begin{aligned}
p(Vu)&=p(u)\quad\text{pour $V$ unitaire dans $F$}\\
p(uU)&=p(u)\quad\text{--- $U$ --- $E$}
\end{aligned}
\]

b) $p(x'\otimes y)=p(x')p(y)$. Supposons $E$ et $F$ de dim infinie pour fixer les idées, \unsure{soient} $(e_i)$ $(f_i)$ des systèmes orthonormaux infinis dans $E$ et $F$, on a une immersion $R^{(\overset{\unsure{\lambda|\varphi}}{\underline N})}\to E'\otimes F\subset L(E,F)$ par $(\lambda_i)\to\sum\lambda_i\bar e_i\otimes f_i$. Alors $p(\varphi(\lambda))$ sur $R^{(N)}$

est une norme, induisant sur $R^n$ des normes \deleted{de Schatten} invariantes par $\unsure{\mathscr G_n|\mathfrak S_n}$ (condition a sur $p$) normalisée (condition b) ne dépendant que de $(|\lambda_i|)$ (condition a)). \deleted{\unsure{On a donc|Elle}} \inserted{$N$ est donc une} norme de Schatten. On vérifie alors immédiatement (condition a)) que l'on a $N(u)=p(u)$ pour \deleted{\uline{tout}} $u\in E'\otimes F$.

On a ainsi obtenu une caractérisation axiomatique des normes $N(u)$ sur les espaces \deleted{$L(E,F)$}. $E'\otimes F$.

\margin{102}
